\chapter{I test di valutazione funzionale nel calcio}
\label{cap:testvalutazionecalcio}

Dal modello prestativo si passa agli \textbf{strumenti di misura} con cui il preparatore raccoglie i
dati su cui costruire la programmazione. Questa parte del corso è divisa in tre lezioni: la prima
--- raccolta qui --- è \textbf{teorica} e fissa che cos'è un test, come si classifica, come si
elaborano i dati raccolti e quali principi reggono la valutazione antropometrica, della flessibilità
e della forza; le successive entrano nel merito dei singoli test usati in ambito calcistico.

% ---------------------------------------------------------------------------
\section{Che cos'è un test}

\textbf{Test}: prova nella quale sono misurate determinate variabili, utili a impostare
l'allenamento; il termine è \textbf{sinonimo di misura}.

\rubrica{A che cosa serve}
\begin{itemize}
  \item è lo \textbf{strumento indispensabile} per la programmazione dei processi di allenamento a
    \textbf{breve, medio e lungo termine};
  \item permette di \textbf{diagnosticare eventuali insufficienze} che non sempre possono essere
    intuite dalla \textbf{semplice osservazione}.
\end{itemize}

% ---------------------------------------------------------------------------
\section{La classificazione dei test}

\subsection{Test generali e test specifici}

\begin{itemize}
  \item \textbf{generali}: effettuati nell'ambito della \textbf{valutazione per categorie};
    indagano le \textbf{caratteristiche fisiche di base} --- forza, potenza, resistenza, velocità,
    mobilità articolare --- e hanno lo scopo di verificare i \textbf{prerequisiti fisici};
  \item \textbf{specifici}: indagano le \textbf{qualità specifiche} e possiedono un'elevata
    \textbf{valenza tecnica}; devono essere strutturati secondo le esigenze della \textbf{disciplina
    sportiva} o del \textbf{singolo atleta}.
\end{itemize}

\subsection{Test da laboratorio e test da campo}

\begin{itemize}
  \item \textbf{da laboratorio}: svolti in \textbf{ambiente controllato}, con protocolli e
    strumentazioni che \textbf{simulano} l'attività sportiva;
  \item \textbf{da campo}: proposti \textbf{mentre l'atleta è impegnato nella propria attività
    sportiva}; sono \textbf{meno affidabili} ma \textbf{più validi}, perché più specifici.
\end{itemize}

In definitiva: laboratorio $\rightarrow$ più \textbf{standardizzabile} ma meno specifico; campo
$\rightarrow$ meno standardizzabile ma \textbf{più specifico}.

\subsection{Test diretti e test indiretti}

\begin{itemize}
  \item \textbf{diretti}: il parametro è misurato direttamente --- p.\,es.\ la massima potenza
    aerobica ($VO_2max$) mediante l'\textbf{analisi dei gas espirati} al termine di un test
    incrementale;
  \item \textbf{indiretti}: il parametro è \textbf{ricavato per correlazione} da una grandezza
    diversa --- p.\,es.\ il \textbf{test di Cooper}, in cui la distanza percorsa in 12$'$ è
    correlata con il $VO_2max$;
  \item nella valutazione funzionale \textbf{sul campo} molti test sono indiretti, perché misurano
    una caratteristica che \textbf{non può essere misurata direttamente} al di fuori del laboratorio.
\end{itemize}

\subsection{Test massimali e test sub-massimali}

\begin{itemize}
  \item \textbf{massimali}: si chiede al giocatore una prova massimale, interrotta al sopraggiungere
    di uno \textbf{stato soggettivo di fatica} incompatibile con la prosecuzione della prova. La
    prestazione può essere:
    \begin{itemize}
      \item \textbf{massimale da subito} --- uno sprint sui 5--10\,m, un test di salto;
      \item raggiunta \textbf{progressivamente} --- le prove incrementali.
    \end{itemize}
  \item \textbf{sub-massimali}: l'intensità proposta \textbf{non è mai massimale} --- p.\,es.\ i
    test di \textbf{soglia anaerobica}, poiché l'intensità corrispondente alle soglie è sempre
    inferiore alla massimale che il giocatore può raggiungere e sostenere.
\end{itemize}

\subsection{Test statici e test dinamici}

\begin{itemize}
  \item \textbf{statici}: \textbf{non richiedono movimento} --- misure antropometriche, frequenza
    cardiaca a riposo, forza isometrica, flessibilità;
  \item \textbf{dinamici}: comportano l'\textbf{attivazione del sistema muscolare} durante movimenti
    specifici, che possono essere:
    \begin{itemize}
      \item in forma \textbf{ciclica} --- la corsa;
      \item in forma \textbf{aciclica} --- il salto.
    \end{itemize}
  \item rientrano fra i dinamici i test di \textbf{forza dinamica} e \textbf{isocinetici}, quelli
    per la determinazione delle \textbf{soglie} e della \textbf{massima potenza aerobica}, e quelli
    di \textbf{agilità}.
\end{itemize}

\subsection{Test singoli e batterie di test}

\begin{itemize}
  \item nella stessa seduta si tende a indagare \textbf{qualità differenti}, somministrando una
    \textbf{batteria di test}: una serie di prove svolte in modo da ricavare più informazioni;
  \item nelle batterie va considerato l'\textbf{intervento delle varie fonti energetiche}: i test
    \textbf{più faticosi} vanno proposti \textbf{alla fine}, dopo adeguato recupero, o addirittura
    in una \textbf{seduta successiva}.
\end{itemize}

% ---------------------------------------------------------------------------
\section{L'elaborazione e l'interpretazione dei dati}

\rubrica{Il criterio generale}
\begin{itemize}
  \item raccolti i dati, il compito è \textbf{analizzare e interpretare} i risultati, per ottenere
    le informazioni con cui \textbf{impostare o modificare} l'allenamento;
  \item il modo migliore di elaborare è \textbf{fare le cose semplici}: si parte dal riportare i
    dati in \textbf{tabelle} debitamente redatte (p.\,es.\ fogli Excel);
  \item i risultati vanno \textbf{discussi} con i tecnici e i colleghi della stessa squadra e con i
    \textbf{singoli giocatori}, che devono essere messi a conoscenza dei propri risultati.
\end{itemize}

\subsection{I dati del singolo giocatore}

\begin{itemize}
  \item compito \textbf{abbastanza semplice}: i dati possono essere raccolti \textbf{come tali} ---
    il tempo impiegato in un test di sprint, il tempo di volo nel test di Bosco;
  \item il passo successivo è l'\textbf{interpretazione} rispetto a \textbf{parametri di
    riferimento}: la scelta dei riferimenti \textbf{appropriati} al test effettuato è fondamentale;
  \item i dati del singolo si confrontano spesso con quelli \textbf{della squadra} o dei giocatori
    che ricoprono lo \textbf{stesso ruolo}; quando è possibile, anche con quelli di \textbf{altre
    squadre}.
\end{itemize}

\subsection{I dati della squadra}

\begin{itemize}
  \item la situazione è \textbf{più complessa} ed è utile qualche \textbf{calcolo statistico};
  \item si costruisce una \textbf{tabella} con i valori del test, che va poi presentata in
    \textbf{tre ordinamenti}:
    \begin{enumerate}
      \item per \textbf{giocatore};
      \item suddivisa per \textbf{ruolo};
      \item per \textbf{prestazione}, dal giocatore con la prestazione migliore a tutti gli altri in
        ordine \textbf{decrescente}.
    \end{enumerate}
\end{itemize}

\subsection{I due livelli di valutazione}

\begin{enumerate}
  \item la \textbf{fotografia della situazione} --- comprendere il significato del quadro rilevato
    per il singolo giocatore o per la squadra;
  \item il \textbf{confronto nel tempo} --- paragonare i risultati dello \textbf{stesso test}
    ripetuto in momenti successivi della stagione, guardando al \textbf{valore assoluto} oppure alla
    \textbf{differenza} tra i test.
\end{enumerate}

% ---------------------------------------------------------------------------
\section{I test antropometrici}

\begin{itemize}
  \item indagano le caratteristiche legate alla \textbf{misura delle dimensioni del corpo umano} che
    possono influenzare la prestazione del calciatore;
  \item se effettuate e interpretate correttamente, permettono di ricavare informazioni su
    \textbf{anatomia}, \textbf{fisiologia}, \textbf{nutrizione} e sulla presenza di eventuali
    \textbf{patologie} o di esiti di patologie del soggetto in esame.
\end{itemize}

\rubrica{Misure antropometriche di routine}
\begin{enumerate}
  \item \textbf{statura};
  \item \textbf{massa corporea} (comunemente detta peso);
  \item \textbf{tessuto adiposo sottocutaneo} --- plicometria;
  \item \textbf{indice di massa corporea} (BMI):
    \[
      BMI\ [kg/m^2] = \frac{\textrm{massa}\ (kg)}{\textrm{altezza}\ (m)^2}
    \]
  \item \textbf{bioimpedenziometria} (BIA), per determinare la \textbf{composizione corporea}:
    \begin{itemize}
      \item si basa sul principio per cui la \textbf{resistenza al passaggio di una corrente
        elettrica} è \textbf{inversamente proporzionale} alla massa grassa del corpo umano;
      \item essendo un'\textbf{indagine medica}, è consigliabile farla effettuare da un
        \textbf{medico esperto} in materia.
    \end{itemize}
\end{enumerate}

% ---------------------------------------------------------------------------
\section{I test di flessibilità e articolarità}

\textbf{Flessibilità}: estensibilità dei \textbf{tessuti peri-articolari}, che permette il movimento
fisiologico dell'articolazione. Il termine si riferisce alla capacità di allungarsi dell'\textbf{unità
muscolo-tendinea}; l'\textbf{articolarità} è invece la possibilità di movimento
\textbf{dell'articolazione}, cioè l'escursione articolare.

\subsection{La terminologia}

\begin{itemize}
  \item \textbf{articolarità}: possibilità di movimento dell'articolazione nell'ambito fisiologico
    dell'escursione articolare;
  \item \textbf{flessibilità}: possibilità di movimento dell'articolazione \textbf{in dipendenza}
    della capacità di allungarsi dell'unità muscolo-tendinea;
  \item \textbf{\textit{range of motion}} (ROM): arco di movimento;
  \item \textbf{lassità}: possibilità di movimento \textbf{anormale} dell'articolazione, dipendente
    da una \textbf{situazione patologica} (p.\,es.\ lesione dei legamenti). In età giovanile è
    frequente, soprattutto a carico dell'articolazione del \textbf{ginocchio};
  \item \textbf{elasticità}: proprietà di un corpo di \textbf{riprendere la forma originaria} dopo
    aver subito una deformazione.
\end{itemize}

\subsection{Il report di valutazione articolare}

Scheda a sei voci, compilata assegnando a ciascun distretto uno dei \textbf{tre livelli} di giudizio.
Per i distretti \textbf{bilaterali} --- ileopsoas, ischiocrurali, retto femorale, adduttori --- il
giudizio è registrato \textbf{separatamente per destra e sinistra}.

\begin{table}[htbp]
  \centering
  \small
  \renewcommand{\arraystretch}{1.3}
  \begin{tabularx}{\textwidth}{|l|X|X|X|}
    \hline
    \textbf{Distretto} & \textbf{Corretto allungamento} & \textbf{Lieve accorciamento} &
      \textbf{Notevole accorciamento} \\
    \hline
    Addominali & \textit{(forza buona)} assume lentamente la stazione seduta da decubito supino &
      \textit{(forza limitata)} il movimento è possibile solo con le braccia estese in avanti &
      \textit{(forza scarsa)} non riesce a sollevare il busto \\
    \hline
    Colonna vertebrale & distanza fronte--rotula da 0 a 10\,cm & distanza tra 10 e 15\,cm &
      distanza $>$ 15\,cm \\
    \hline
    Ischiocrurali & con un arto fissato si solleva l'arto esteso a 90$^\circ$ senza dolori nella
      zona poplitea & l'arto si solleva con angolo di flessione dell'anca da 80$^\circ$ a 90$^\circ$ &
      l'arto può essere sollevato da 60$^\circ$ a 80$^\circ$ \\
    \hline
    Retto femorale & \textit{(allungamento ottimale)} il tallone raggiunge il gluteo con un leggero
      aiuto passivo & distanza tallone--gluteo di circa 15\,cm & distanza $>$ 15\,cm \\
    \hline
    Adduttori & \textit{(corretta estensibilità)} adduzione di 60$^\circ$ & adduzione da 40$^\circ$
      a 60$^\circ$ & adduzione da 25$^\circ$ a 40$^\circ$ \\
    \hline
  \end{tabularx}
\end{table}

\rubrica{Ileopsoas}
Sesta voce della scheda, senza soglie numeriche: l'\textbf{escursione di movimento della coscia}
dipende dal \textbf{grado di accorciamento dell'ileopsoas controlaterale}. Nei calciatori si nota
spesso un'importante \textbf{retrazione} di questo muscolo.

\rubrica{Perché questi distretti}
\begin{itemize}
  \item gli \textbf{ischiocrurali} sono un distretto soggetto a \textbf{continui infortuni};
  \item l'\textbf{articolarità della caviglia} è molto importante nel calciatore;
  \item la \textbf{mobilità scapolo-omerale} riguarda in modo prioritario i \textbf{portieri}, ma
    non va sottovalutata come dato per la fotografia dello \textit{status} funzionale dell'atleta;
  \item l'\textbf{equilibrio} va valutato sia a \textbf{occhi aperti} sia a \textbf{occhi chiusi}.
\end{itemize}

\subsection{I valori di riferimento}

I valori normativi della scheda sono ripresi da \textbf{bibliografia inglese}. Le prove sono lo
\textit{Static Ankle Test}, lo \textit{Static Shoulder Test} e lo \textit{Standing Stork Test}
(equilibrio a occhi aperti), quest'ultimo replicato a occhi chiusi come \textit{Standing Stork
Blind}.

\begin{table}[htbp]
  \centering
  \small
  \renewcommand{\arraystretch}{1.3}
  \begin{tabularx}{\textwidth}{|X|c|c|c|c|c|c|}
    \hline
    & \multicolumn{2}{c|}{\textbf{Ankle test}} & \multicolumn{2}{c|}{\textbf{Shoulder test}} &
      \multicolumn{2}{c|}{\textbf{Stork test}} \\
    \cline{2-7}
    \textbf{Livello} & M & F & M & F & M & F \\
    \hline
    \textit{Excellent} & $>$89 & $>$81 & $<$17 & $<$12 & $>$50$''$ & $>$30$''$ \\
    \hline
    \textit{Good} / \textit{above average} & 89--82,5 & 81--77,5 & 29--17 & 24--12 & 50--41 &
      30--23 \\
    \hline
    \textit{Average} & 82,5--75 & 77,5--67 & 36--29 & 33--24 & 40--31 & 22--16 \\
    \hline
    \textit{Fair} / \textit{below average} & 75--67 & 67--61,5 & 50--36 & 45--33 & 30--20 &
      15--10 \\
    \hline
    \textit{Poor} & $<$67 & $<$61,5 & $<$49,5 & $<$45 & $<$20$''$ & $<$10$''$ \\
    \hline
  \end{tabularx}
\end{table}

Lo \textit{Standing Stork Blind} non è classificato per livelli ma convertito in un
\textbf{punteggio} a partire dal \textbf{miglior tempo} tenuto:

\begin{table}[htbp]
  \centering
  \small
  \renewcommand{\arraystretch}{1.2}
  \begin{tabularx}{\textwidth}{|X|c|c|c|c|c|c|c|c|c|c|c|c|}
    \hline
    \textbf{\textit{Best time}} & 60$''$ & 55$''$ & 50$''$ & 45$''$ & 40$''$ & 35$''$ & 30$''$ &
      25$''$ & 20$''$ & 15$''$ & 10$''$ & 5$''$ \\
    \hline
    M & 20 & 18 & 16 & 14 & 12 & 10 & 8 & 6 & 4 & 3 & 2 & 1 \\
    \hline
    F & --- & --- & --- & --- & --- & 20 & 17 & 14 & 11 & 8 & 4 & 2 \\
    \hline
  \end{tabularx}
\end{table}

\rubrica{Le soglie dell'equilibrio nell'esperienza del docente}
Riferite a entrambe le condizioni (occhi aperti e occhi chiusi) e rilevate separatamente per gamba
destra e sinistra:
\begin{itemize}
  \item \textbf{45$''$}: valore massimo raggiungibile $\rightarrow$ valutazione \textbf{ottimale};
  \item \textbf{sotto i 30$''$}: \textbf{sufficienza};
  \item \textbf{sotto i 20--15$''$}: \textbf{scarsità} del parametro di equilibrio.
\end{itemize}

% ---------------------------------------------------------------------------
\section{I test di forza e di potenza muscolare}

La \textbf{forza muscolare} è una componente indispensabile per costruire la prestazione di chi
pratica sport.

\rubrica{A che cosa servono i test di forza}
\begin{enumerate}
  \item determinare i \textbf{carichi specifici} da utilizzare per l'allenamento;
  \item \textbf{monitorare gli effetti} dell'allenamento;
  \item \textbf{monitorare il percorso riabilitativo} del giocatore infortunato.
\end{enumerate}

\subsection{I fattori anatomici che determinano la forza}

\begin{itemize}
  \item \textbf{superficie della sezione muscolare}: tanto maggiore è la superficie di sezione,
    tanto maggiore sarà la forza prodotta in qualsiasi situazione;
  \item \textbf{architettura muscolare}: la disposizione delle fibre nella loro inserzione sul
    tendine determina una diversa capacità di produrre forza $\rightarrow$ i muscoli con fibre
    inserite \textbf{obliquamente} al tendine (\textbf{fibre pennate}, p.\,es.\ il deltoide)
    esprimono forze maggiori dei muscoli \textbf{fusiformi} (p.\,es.\ il bicipite brachiale);
  \item \textbf{tipologia muscolare}: \textbf{geneticamente determinata}, pur essendo in qualche
    modo influenzabile dall'allenamento.
\end{itemize}

\rubrica{Classificazione di Ranvier}
\begin{enumerate}
  \item \textbf{ST} (\textit{slow twitch}): fibre a \textbf{lenta contrazione}, caratterizzate dal
    metabolismo \textbf{ossidativo} (aerobico);
  \item \textbf{FT} (\textit{fast twitch}): fibre a \textbf{rapida contrazione}, con metabolismo
    \textbf{anaerobico}, soggette facilmente all'affaticamento;
  \item le FT si suddividono a loro volta in:
    \begin{itemize}
      \item \textbf{FT IIA}: possono usare anche il metabolismo aerobico $\rightarrow$ ``fibre
        rapide e resistenti'' (\textbf{FR});
      \item \textbf{FT IIB}: essenzialmente anaerobiche $\rightarrow$ ``fibre rapide che si
        affaticano'' (\textbf{FF}).
    \end{itemize}
\end{enumerate}

\subsection{I fattori fisiologici che determinano la forza}

\rubrica{Velocità di contrazione}
\begin{itemize}
  \item tanto maggiore è la velocità di contrazione, tanto \textbf{minore} sarà la forza espressa;
  \item la diminuzione della forza massima all'aumentare della velocità dipende dal \textbf{tempo
    necessario a formare tutti i ponti acto-miosinici}, che decresce all'incrementare della
    velocità;
  \item la \textbf{relazione forza/velocità} è fondamentale perché permette di classificare la forza
    muscolare sulla base dell'\textbf{azione muscolare}, distinta in \textbf{statica} e
    \textbf{dinamica}.
\end{itemize}

\rubrica{Fattori metabolici}
\begin{itemize}
  \item in particolare la \textbf{disponibilità dei substrati energetici}, che influenza la capacità
    di \textbf{prolungare} e/o \textbf{ripetere} l'esercizio nel tempo;
  \item ruolo centrale nella \textbf{genesi della fatica}, che si traduce in una progressiva
    diminuzione della forza muscolare;
  \item la fatica \textbf{non è solo un fenomeno acuto}: per le esercitazioni che coinvolgono il
    \textbf{ciclo stiramento-accorciamento} la capacità di produrre forza può risultare depressa
    anche per \textbf{alcuni giorni}, con ovvie ricadute sulla \textit{performance}.
\end{itemize}

\rubrica{Fattori nervosi}
Il muscolo scheletrico si contrae in risposta a un comando che origina volontariamente dal
\textbf{sistema nervoso centrale}; i fattori nervosi sono numerosi e molto complessi:
\begin{enumerate}
  \item \textbf{soglia di eccitabilità} dei motoneuroni;
  \item \textbf{velocità di conduzione nervosa};
  \item \textbf{numero di fibre} muscolari contratte \textbf{simultaneamente};
  \item efficienza della \textbf{sincronizzazione};
  \item grado di \textbf{inibizione} delle unità motorie non attivate.
\end{enumerate}

% ---------------------------------------------------------------------------
\section{La valutazione della forza}

Vari autori hanno sottolineato l'importanza della forza nel calcio; tuttavia \textbf{scegliere quali
test adottare} per la valutazione funzionale del calciatore può essere un compito arduo per il
preparatore atletico e per lo staff tecnico.

\subsection{Le due situazioni biomeccaniche}

Effettuando un test di forza si è \textbf{sempre} in una di queste due condizioni:

\begin{enumerate}
  \item \textbf{si conosce la forza richiesta}, che corrisponde al \textbf{carico} spostato da un
    determinato gruppo muscolare (p.\,es.\ 10\,kg), oppure al \textbf{peso corporeo} del calciatore
    (p.\,es.\ in un salto) $\rightarrow$ essendo noto il carico, va misurata la \textbf{velocità} con
    cui viene spostato;
  \item \textbf{si conosce la velocità} del movimento, perché \textbf{preimpostata} dagli apparecchi
    \textbf{isocinetici} $\rightarrow$ va misurata la \textbf{forza espressa}. È il caso anche dei
    più comuni \textbf{dinamometri isometrici}, dove la velocità di movimento è \textbf{nulla} e
    costante; la valutazione isometrica è però considerata \textbf{aspecifica} rispetto al gesto
    sportivo.
\end{enumerate}

\subsection{La specificità}

Vanno sempre considerate le \textbf{caratteristiche specifiche dei movimenti} impiegati per i test,
sulla base delle informazioni che si desidera ricavare, distinguendo i test che indagano le
caratteristiche \textbf{generali} da quelli che indagano le \textbf{specifiche} --- che dovrebbero
essere il più possibile \textbf{simili a quanto avviene in gara}.

\rubrica{I sette aspetti da considerare}
\begin{enumerate}
  \item \textbf{complessità dei movimenti}: mono- o pluriarticolari;
  \item \textbf{postura}: influisce sull'attivazione e/o sull'inibizione di determinati gruppi
    muscolari;
  \item \textbf{arco di movimento}: la forza varia durante il movimento articolare;
  \item \textbf{tipo di azione muscolare}: isometrica, concentrica, eccentrica, o che coinvolga il
    ciclo stiramento-accorciamento;
  \item \textbf{entità della forza}: forza massima, forza media, momento di forza, oppure le
    \textbf{forze di reazione con il terreno} --- aspetto tutt'altro che secondario, perché il
    calciatore gioca su fondi molto diversi (il terreno d'inverno ha una consistenza differente, e
    su sintetico la forza di reazione è un'altra);
  \item \textbf{parametri cinematici}: velocità e accelerazione, che a loro volta variano durante i
    movimenti;
  \item \textbf{tipo di apparecchiatura} utilizzata: macchine o pesi liberi.
\end{enumerate}

\subsection{Il quadro riassuntivo dei test di forza}

\begin{table}[htbp]
  \centering
  \small
  \renewcommand{\arraystretch}{1.3}
  \begin{tabularx}{\textwidth}{|l|X|X|}
    \hline
    \textbf{Tipo di test} & \textbf{Caratteristiche} & \textbf{Variabili misurabili} \\
    \hline
    Forza isometrica & le condizioni statiche annullano la variabile velocità di movimento e il
      muscolo esprime la sua forza massima; la forza dipende dall'\textbf{angolo articolare}, cioè
      dai vantaggi meccanici e dalla lunghezza del muscolo & forza massima isometrica \\
    \hline
    Forza isocinetica & si impone una \textbf{velocità costante} di movimento e si misura la forza
      dinamica; la forza varia durante tutto l'arco di movimento in base all'angolo articolare &
      momento di forza e \textbf{picco} del momento di forza; momento di forza a un determinato
      angolo articolare \\
    \hline
    Forza massima & si usano i \textbf{sovraccarichi} per misurare la forza massima di un
      determinato gruppo muscolare & 1 RM \\
    \hline
    Potenza & movimenti eseguiti alla \textbf{massima velocità possibile}, tipo salto in alto o
      salto in lungo & durata del movimento; distanza raggiunta; potenza espressa \\
    \hline
  \end{tabularx}
\end{table}
